%=====================================================================
%  ch04 —— 原文 863–1207：§3 Beilinson–Bernstein 局部化（3.1 概览、3.2 局部化定理）
%=====================================================================

\section{Beilinson--Bernstein 局部化与隐藏的对称性}
\label{s:bbloc}

\subsection{概览}
\label{s:overview3}

\ssszh 上一节我们考虑的是旗流形的 $\Fq$-点上函数构成的向量空间
%
\begin{equation} \label{e:funflgs}
\Fun(\sB(\Fq)),
\end{equation}
%
它自然地带有 $G(\Fq)$ 的作用。

本节我们改为关注复旗流形上 D-模的导出范畴
%
\begin{equation} \label{e:dmodflg}
\Dmod(\sB(\C)).
\end{equation}
%
我们先按 Beilinson--Bernstein 回顾这个范畴 \eqref{e:dmodflg} 与李代数表示论之间的基本
联系。在回顾这一对应的若干重要性质时，我们将被引向 $G(\Fq)$ 在 $\Fun(\Fl_G(\Fq))$ 上
作用的类似物。由此我们就会遇到群的范畴作用这一想法。

\subsection{局部化定理}

\ssszh 读者或许已从第 \ref{s:overview3} 小节的讨论中看出，人们把 D-模范畴
\eqref{e:dmodflg} 看作 \eqref{e:funflgs} 的某种范畴化。这个类比在几何表示论中是绝对
根本的。为此，我们在附录 \ref{s:frmfun} 中对它作了相当详细的解释，并介绍了 D-模。
非专家读者不妨先翻到那里，掌握了基础再回来读本节。

\ssszh 本节余下部分在特征零的代数闭域 $k$ 上工作。为记号简便，我们重新定义 $\Fl_G$ 为
相应的基变换。也就是说，若用 $\Fl_{G, \mathbb{Z}}$ 记第 \ref{s:repthryfingroup} 节
讨论过的 $\Spec \Z$ 上的版本，则置
%
\[
 \Fl_G := \Fl_{G, \Z} \underset {\Spec \Z} \times \Spec k.
\]
%
当然，读者可以取 $k$ 为复数，但要注意我们并不需要它们的度量拓扑的任何特殊性质。

\ssszh 旗流形上的 D-模范畴
$$\Dmod(\Fl_G)$$
原来还有一副看似不同的面貌：它可以实现为李代数表示的范畴。下面回顾这一孕育了几何表示论
的非凡对应。

\ssszh 回忆光滑代数簇 $X$ 上的 D-模 $\sM$ 是带平坦联络的拟凝聚层。等价地，$\sM$ 带有
微分算子层 $\sD_X$ 的作用，且该作用延拓正则函数 $\OO_X$ 的作用。特别地，其整体截面带有
整体微分算子的作用，即有一个重言函子
%
\begin{equation} \label{e:glob}
\Gamma(X, -): \Dmod(X) \rightarrow \Gamma(X, \sD_X)\on{-mod.}
\end{equation}
%
Beilinson--Bernstein 的一个著名定理说：当 $X = \sB$ 时，这个函子是等价。

\ssszh 我们来说明这为什么出人意料。

首先，若 $X$ 是仿射的，类似的断言成立，但并不出人意料。事实上，此时它来自拟凝聚层
定义的设定方式：在仿射簇上，函子
%
\begin{equation} \label{e:globqcoh}
\Gamma(X, -): \on{QCoh}(X) \rightarrow \Gamma(X, \OO_X)\mod
\end{equation}
%
是等价。

相反，若 $X$ 不是仿射的，例如射影的，则 \eqref{e:globqcoh} 几乎从不是等价。例如，
对 $\sB$ 而言，其整体（导出）函数只有标量，于是 \eqref{e:globqcoh} 退化为映射
%
\[
\Gamma(\sB, -): \on{QCoh}(\sB) \rightarrow \on{Vect,}
\]
%
其中 $\on{Vect}$ 记向量空间的导出范畴。这远非等价，因为 $\sB$ 不是一个点！正因如此，
在进一步给拟凝聚层配上平坦联络之后，\eqref{e:glob} 对旗流形成立为等价，才是非平凡的事。

我们还要强调：对一般射影簇，\eqref{e:glob} 不是等价。例如，熟悉代数曲线的读者可以
验证：使 \eqref{e:glob} 成为等价的光滑射影曲线只有 $\mathbb{P}^1$，即最小的旗流形。

\ssszh 下面转向整体微分算子
%
\[
\Gamma(\sB, \sD_{\sB})
\]
%
的描述。回忆对任意光滑簇 $X$，$\sD_{X}$ 作为代数层由函数 $\OO_X$ 与向量场 $\sT_X$
生成。我们来看在本情形 $X = \sB$ 中，它们对整体截面有何贡献。

如上所述，旗流形上的整体函数只是标量
%
\[
k \xrightarrow{\sim} \Gamma(\sB, \OO_{\sB}).
\]
%
对向量场，$G$ 在 $\sB$ 上的自然作用给出其李代数 $\fg$ 的无穷小作用，即李代数同态
%
\begin{equation}\label{e:vecfields}
\fg \rightarrow \Gamma(\sB, \sT_{\sB}).
\end{equation}
%
为接近最终答案，我们可以乐观地希望这是同构。

为顾及高阶微分算子，我们同样可以要求由包络代数出发的映射
%
\begin{equation}\label{e:globdiff}
U(\fg) \rightarrow \Gamma(\sB, \sD_{\sB})
\end{equation}
%
是同构。这将是「$\sD_{\sB}$ 的截面局部地由函数与向量场生成」的一个强整体类比。

\ssszh 结果，\eqref{e:vecfields} 与 \eqref{e:globdiff} 都不是同构。必须考虑中心来修正。

先从向量场说起。注意 $G$ 的中心 $Z$ 平凡地作用在 $\sB$ 上。因此需要模去 $\fg$ 的中心
$\fz$。这样做之后，确实得到等价
%
\begin{align*}
\fg/\fz &\xrightarrow{\sim} \Gamma(\sB, \sT \hspace{.2mm}).
\intertext{类似地，过渡到所有微分算子时，必须施加中心商。也就是说，若用 $Z(\fg)$
记 $U(\fg)$ 的中心，它在 $\fg$ 的平凡表示上的作用给出一个特征
$\chi: Z(\fg) \rightarrow k.$ 用 $U_0(\fg)$ 记整个包络代数相应的中心商}
U_0(\fg) &:= U(\fg) \underset{Z(\fg)} \otimes  k_\chi.
\intertext{于是与向量场的情形一样，\eqref{e:globdiff} 穿过所得的商，并得到等价}
U_0(\fg) &\xrightarrow{\sim} \Gamma(\sB, \sD).
\end{align*}

\label{s:casa}

\begin{ex} 我们显式描述 $\fg = \fgl_2$ 时 $U_0(\fg)$ 的样子。可把李代数分解为中心
与无迹矩阵
%
\[
\fgl_2 = \fz \oplus \fsl_2.
\]
%
这诱导包络代数及其中心的分解
%
\[
U(\fgl_2) \simeq \on{Sym}(\fz) \otimes U(\fsl_2) \quad \text{与} \quad
Z(\fgl_2) \simeq \on{Sym}(\fz) \otimes Z(\fsl_2).
\]
%
\label{s:cas}
中心 $Z(\fsl_2)$ 是单变量多项式环，由滤过次数为二的 Casimir 元 $\Omega$ 生成。
具体地说，若取 $\fz$ 的基元 $z$ 与 $\fsl_2$ 的标准生成元 $f, h, e$，则
%
\[
Z(\fgl_2) \simeq k[ z, ef + fe + \frac{1}{2}h^2],
\]
%
其中第二个生成元就是 $\Omega$。容易看出两个生成元在平凡表示上都以零作用。特别地，
$U_0(\fg)$ 是 $U(\fgl_2)$ 模去这两个元素所生成的理想得到的商，即
%
\[
U_0(\fgl_2) = U(\fgl_2)/( \fz, \Omega) \simeq U(\fsl_2)/(\Omega).
\]
%
\end{ex}

\ssszh 用 $\fg\on{-mod}_0$ 记 $U_0(\fg)$-模的导出范畴，即中心特征平凡的 $\fg$ 表示。
总结上面的讨论，Beilinson--Bernstein 定理如下。

\begin{theo} \label{t:bbloc}(Beilinson--Bernstein \cite{beilinson-bernstein81})
整体截面函子给出 $t$-正合的等价
%
\[
\Gamma(\sB, -): \Dmod(\sB) \rightarrow \fg\on{-mod}_0.
\]
%
\end{theo}

这个定理对表示论发展的意义，怎么强调都不过分。这里只作几点说明。

\ssszh 第一，由 Quillen 的一个 Schur 引理版本，每个单 $\fg$-模都有中心特征，即中心
$Z(\fg)$ 以标量作用。因此，研究不可约模时，逐个中心特征来做就足够了。

上述定理把中心特征平凡的那些模，与 $\sB$ 上的单 D-模等同起来。这使人得以把这类不可约模
研究中的问题，翻译成几何或拓扑的问题。后者可以借助代数簇拓扑的丰富知识来处理，
特别是 Hodge 理论及其奇异变体。

该定理最初的动机正是这样一种应用。即在所有 $\fg$-模的范畴中，对 $G$ 的任意 Borel
子群 $B$，可联系上 $(\fg, B)$-模范畴，即其 $B$ 的李代数作用可积的 $\fg$-模。

一个自 Verma 1966 年的论文 \cite{verma66} 起受到大量研究的问题，是确定单
$(\fg, B)$-模的特征。在 Jantzen \cite{jantz79} 等人于若干低秩情形定出之后，Kazhdan 与
Lusztig 于 1979 年对中心特征平凡的单 $(\fg,B)$-模提出了一般猜想 \cite{kl79}。提出猜想
时，他们已知道其公式与 $\sB$ 的 Schubert 子簇密切相关，更确切地说，与其相交上同调
D-模的 $!$-茎密切相关。但与 $\fg$-模的联系尚不清楚。

Beilinson--Bernstein 的工作 \cite{bb82}，以及独立的 Brylinski--Kashiwara
\cite{brylin81}，提供了所需的联系。即局部化定理把 $(\fg,B)\mod_0$ 的单对象与上述相交
上同调层互换，特征公式随之而来。

\begin{re} 上面只讨论了平凡中心特征的情形。Beilinson--Bernstein 实际上对任意正则中心
特征证明了类似的陈述。关于包括广义或奇异中心特征的变体，见 \cite{kashloc}、
\cite{beilinsonginzburg}、\cite{bmr}、\cite{krem}、\cite{ahc}。 \end{re}

\ssszh 接下来我们用三个基本例子说明局部化等价的行为。

\begin{ex} 如第 \ref{s:casa} 小节所述，局部化把微分算子代数与包络代数的中心商互换，即
%
\[
\Gamma(\sB, \sD) \simeq U_0(\fg).
\]
%
\end{ex}

\begin{ex}如我们已讨论的，局部化把结构层与平凡表示互换，即
%
\[
\Gamma(\sB, \OO) \simeq k.
\]
%
读者或许知道，这是 Borel--Weil--Bott 定理的最简单情形；该定理把所有不可约 $G$-模构造为
$\sB$ 上线丛的整体截面。非正式地说，Beilinson--Bernstein 的工作把这一点推广，给出
所有不可约 $\fg$-模的类似构造。
\end{ex}

\begin{ex} 最后，对 $\sB$ 的点 $x$，考虑 $\delta_x$ 的像。回忆 $x$ 的稳定化子是 Borel
子群 $B$。于是，用 $\fb$ 记 $B$ 的李代数，delta D-模被送到相应的诱导模
%
\begin{equation}
\Gamma(\Fl_G, \delta_x) \simeq U(\fg) \underset{U(\fb)} \otimes
\det(\fg/\fb).
\end{equation}
%
我们略述其证明，读者不妨自己完整想一遍。为造出映射
%
\begin{equation} \label{e:hwv}
U(\fg) \underset{U(\fb)} \otimes k_{\chi} \rightarrow \Gamma(\Fl_G,
\delta_x),
\end{equation}
%
注意把右边的生成线与 $\det(T_x \sB) = \det( \fg/\fb)$ 等同起来，与 $\fb$-作用是相容的，
参看注 \ref{r:hwdelta}。要看出所得的映射是同构，只需注意两边都由任一与 $\fb$ 横截的
子代数 $\fn^-$（即某个对立 Borel 的幂幺根）的作用，从各自的行列式线自由生成。
\end{ex}
